\chapter{Background and related work}
\label{cha:background}

The design defended in this thesis combines two research traditions that have each solved half of the dietary planning problem. Operations research settled the arithmetic side in the 1940s, producing diets that are provably optimal but often unappetising. Natural language processing has more recently produced systems that interpret what a person asks for and can write a recipe they might enjoy, but that cannot be trusted to add up the recipe's contents. This chapter reviews both traditions and the machinery between them: the failure modes that make an unaided language model unsuitable for this task, the tool-use and self-correction techniques that mitigate those failures, and the evaluation methods that make claims about such systems measurable. It closes with the recent dietary systems this work is most directly comparable to and the gap it sets out to fill.

\section{Diet planning as a mathematical optimisation problem}
\label{sec:background-diet-optimisation}

Dietary planning was recognised as a problem with a precise mathematical structure long before language models existed. In 1945, Stigler posed the cheapest combination of \num{77} foods that would satisfy the recommended daily allowances of nine nutrients for a moderately active man, solving it heuristically~\cite{stigler1945cost}; Dantzig later reported that this heuristic came within \num{24} cents a year of the true optimum~\cite{dantzig1990diet}. Dantzig reconstructed the diet problem as an early test case for the simplex algorithm. Jack Laderman of the National Bureau of Standards Mathematical Tables Project computed the optimum -- \textdollar\num{39.69} a year against Stigler's heuristic \textdollar\num{39.93} -- and Dantzig's account of that computation established the \emph{diet problem} as a standard illustration of linear programming~\cite{dantzig1990diet}. The formulation is simple: nutrient requirements become linear inequality constraints, food quantities become continuous decision variables, and a cost or deviation term becomes the objective to minimise.

Subsequent work refined the formulation, not the underlying idea. Sklan and Dariel extended it to mixed-integer linear programming (MILP) so that whole-unit foods and higher-level constraints such as meal structure or exchange-group frequency could be represented directly~\cite{sklan1993diet}. A 2018 review of fifty-two such studies concludes that most imposed nutritional and cost constraints, that only six offered a good example of a palatability or acceptability constraint, and that none of the reviewed studies offered a fully satisfactory way to encode how acceptable an optimal diet is to the person who has to eat it~\cite{vandooren2018review}. This is the central limitation of the classical approach: a solver finds the numerically optimal diet with respect to its constraint matrix, but has no notion of a recipe, a cuisine, or a craving, and encoding those as linear constraints only approximates something inherently linguistic and personal.

This gap is where large language models are attractive: they excel at the semantic, preference-laden reasoning a solver cannot express, while a solver excels at the arithmetic a language model cannot be trusted to perform. A 2026 preprint combining GPT-4o with a MILP solver over a shared \num{297}-item food catalogue makes the trade-off explicit; a team of five evaluators -- medical doctors, public-health specialists and dietitians -- rated one meal per approach for each of five synthetic patient profiles on a five-point scale. On nutrient accuracy the solver alone scored \num{4.93} and the language model alone \num{1.54}, with the hybrid between them at \num{3.96}. The hybrid led on personalisation (\num{3.81} against the solver's \num{3.43} and the model's \num{2.60}) and was rated significantly more practical than the solver ($p = \num{0.021}$); the solver still led on variety (\num{4.10})~\cite{alkeyeva2026chatgpt}. This thesis splits the same labour at a different point and without a general-purpose solver: the model chooses the foods and fixes their portions, leaving the deterministic side no free variables to optimise over and only a sum to compute (\Cref{sec:design-totaller}).

\section{Failure modes of large language models in numeric and safety-critical tasks}
\label{sec:background-llm-failure-modes}

The transformer architecture~\cite{vaswani2017attention} underlying contemporary large language models is trained, in those models, to predict the next token given its predecessors~\cite{brown2020language}. Nothing in this objective rewards tracking a running total correctly, refusing to answer when a fact is unknown, or applying a hard constraint consistently across a long output; each is at best an emergent side effect of pretraining. Three failure modes that follow are directly relevant to a dietary recommendation system.

The first is \emph{hallucination}: producing fluent content that is unfaithful to, or unsupported by, any real source, including nutrient values for foods the model has never seen data for~\cite{ji2023survey}. A hallucinated number is indistinguishable from a correct one in the output alone, so the only reliable defence is to never let the model be the source of a number that can instead be looked up or computed. The second is unreliable multi-step arithmetic: even given every ingredient's nutrient values, summing them across a full day of several meals is a long chain of deterministic steps of the kind transformer models have been shown to execute unreliably as compositional complexity grows~\cite{dziri2023faith}. The third is \emph{sycophancy}: deferring to a user's stated preference even when it conflicts with an earlier, medically motivated constraint. In an evaluation of three commercial models subjected to rebuttals of escalating strength across mathematical and medical question answering, \qty{14.66}{\percent} of responses were \emph{regressively} sycophantic: the model abandoned a correct answer for the user's incorrect one. Rebuttals appending an authoritative-sounding citation and abstract produced the highest regressive rate. Sycophancy of either kind also proved persistent: once a model had conceded, it stayed conceded through the rest of the rebuttal chain, with at most one transition, in \qty{78.5}{\percent} of chains~\cite{fanous2025syceval}. For a system that must keep a user's sodium intake low regardless of how insistently they ask for something salty, sycophancy is as serious a hazard as hallucination.

These failure modes motivate the recurring architectural principle of this thesis: whatever can be checked mechanically should be, and whatever the model outputs that cannot be checked should be scoped as narrowly as possible.

\section{Knowledge-based and neuro-symbolic systems}
\label{sec:background-neuro-symbolic}

A \emph{knowledge-based system} grounds its behaviour in an explicit, inspectable body of facts and rules, not solely in statistical patterns learned from data~\cite{hayesroth1983knowledge}. Classical dietary knowledge-based systems -- the programmes of \Cref{sec:background-diet-optimisation} -- are entirely symbolic: their inputs, constraints and outputs are explicit and auditable, at the cost of being unable to converse or interpret an ambiguous natural-language request. \emph{Neuro-symbolic} systems combine such symbolic components with a neural one, using the neural component for flexible pattern recognition or language understanding and the symbolic component for exact or auditable reasoning~\cite{garcez2023neurosymbolic}.

The architecture in \Cref{cha:design} follows this pattern at the level of individual data items, not that of a general theorem prover: the language model produces a structured plan that references food knowledge base entries, and a separate deterministic component sums those entries independently of anything the model claims.

\section{Tool-augmented language model agents}
\label{sec:background-tool-use}

For a language model to consult an external, ground-truth source instead of its own parametric memory, it must decide when to call an external function, with what arguments, and how to fold the result back into its reasoning. Two lines of work established this pattern.

Prompting a sufficiently large model to produce an explicit \emph{chain of thought} before its answer improves multi-step reasoning without changing the model's weights~\cite{wei2022chain}, but such a chain still draws only on the model's parametric memory. \textsc{ReAct} interleaves free-text reasoning traces with concrete actions issued to an external environment, feeding each action's observation back into context before the next step. This lets a model revise its plan in light of new information -- retrying a lookup under a different name when the first returns nothing. On HotpotQA, in a hand-labelled sample of fifty failing trajectories per method, the share of failures attributed to hallucination fell from \qty{56}{\percent} under chain-of-thought to none under \textsc{ReAct}, even though exact-match accuracy on the full set was slightly lower (\qty{27.4}{\percent} against \qty{29.4}{\percent})~\cite{yao2023react}.

\textsc{Toolformer} pursues the same goal from the training side, teaching a model in a self-supervised manner to decide which tool to invoke, with which arguments, and where to splice in the result, without a human-annotated training set, bootstrapping from a handful of in-prompt demonstrations per API~\cite{schick2023toolformer}. The relevant lesson is architectural: a model need not \emph{know} a nutrient value or a running total from its own weights if it is given a well-specified interface to retrieve or compute it.

A related mechanism is \emph{structured output}: constraining a model's response to a predeclared schema instead of free text. This does not prevent a hallucinated value from appearing in a well-typed field, but it turns the response into a data structure a program can inspect and reject before anything downstream consumes it. What this thesis builds on that is taken up in \Cref{sec:tools-agent-framework}.

\section{Self-correction: critique-and-refine loops}
\label{sec:background-self-correction}

A complementary mechanism for reducing erroneous outputs is to have a system review and revise its own draft before returning it.

\textsc{Self-Refine} formalises this as an iterative loop in which the model that produced an output critiques it against the task requirements and then revises it in light of that critique, repeating until no further issues are reported or a budget is exhausted~\cite{madaan2023selfrefine}. Because critique and revision come from the same ungrounded model, the approach is easy to implement but offers no guarantee the critique is more reliable than the patterns that produced the original error.

\textsc{CRITIC} closes this gap by grounding the critique in external tool feedback: a draft is checked against a search engine, code interpreter, or other verification tool, and only the externally verified feedback drives revision~\cite{gou2024critic}. Across question answering, program synthesis and toxicity reduction, tool-grounded critique produced larger gains than self-reflection alone. A related study of hallucination in medical question answering formalises the same generate-score-refine pattern with explicit thresholds on entailment, factuality and consistency scores that a candidate must clear before the loop terminates~\cite{ji2023towards}.

Where this thesis sits between those designs -- a tool-less critic still shown the totaller's computed totals -- is taken up in \Cref{sec:tools-self-correction}.

\section{Evaluating generative systems: LLM-as-a-judge}
\label{sec:background-evaluation}

Reference-based metrics -- \textsc{BLEU}~\cite{papineni2002bleu} and \textsc{ROUGE}~\cite{lin2004rouge} -- score a candidate by its n-gram overlap with fixed reference texts and correlate poorly with human judgement on tasks rewarding creativity over surface overlap. On SummEval, ROUGE-L reaches a summary-level Spearman correlation of only \num{0.165} with human ratings when averaged across the four quality dimensions, and on Topical-Chat the same metrics reach turn-level correlations of \num{0.244} (ROUGE-L) and \num{0.259} (BLEU-4) averaged across the four aspects~\cite{liu2023geval}. A meal plan's rationale and recipe instructions are exactly this kind of output: there is no single correct phrasing, and two equally good plans may share almost no vocabulary. \textsc{G-Eval} prompts a strong language model to generate a chain-of-thought evaluation procedure from a natural-language description of the criterion, then apply it to the candidate in a form-filling manner, reporting a probability-weighted score; on SummEval this raised the summary-level Spearman correlation with human ratings well above those overlap baselines~\cite{liu2023geval}. A broader study using GPT-4 to judge pairs of chatbot responses on MT-Bench found that the judge matched expert pairwise preferences on \qty{85}{\percent} of non-tied comparisons, against \qty{81}{\percent} agreement between two human raters, while documenting systematic biases towards longer responses and towards position in a pairwise comparison~\cite{zheng2023judging}. The qualitative scoring stage of this thesis follows the G-Eval pattern; its departures -- hand-written steps per criterion, a Monte Carlo estimate of the probability-weighted score, and two judges -- are specified in \Cref{sec:methodology-qualitative,sec:methodology-judge-protocol}.

\section{Related dietary recommendation systems}
\label{sec:background-related-systems}

Two recent systems combine a language model with an external knowledge source for dietary recommendation in ways directly comparable to this thesis. \textsc{NutriGen} relies on prompt engineering, not an output schema or deterministic totaller, to incorporate USDA reference values into a generated meal plan; its evaluation reports a mean caloric-target error of \qty{1.55}{\percent} for Llama~3.1~8B against \qty{8.08}{\percent} for the larger Llama~3.1~70B, and \qty{3.68}{\percent} for GPT-3.5~Turbo against \qty{13.47}{\percent} for GPT-4o, showing that raw model capability does not reliably predict numeric accuracy here~\cite{khamesian2025nutrigen}. \textsc{DietQA} takes the opposite route and answers from a corpus of dishes it has already collected: a crawler harvests recipe sites, each ingredient line is mapped to a canonical entry whose nutrient profile is imported from a reference database, and the result is stored as a knowledge graph carrying dietary tags. Hard exclusions (an allergen, a diet pattern) are then enforced by graph queries before the model narrates the recipes that survive~\cite{tsampos2025dietqa}.

Both demonstrate that pairing a language model with a queryable knowledge source measurably constrains its output, and \textsc{DietQA} in particular already computes nutrient totals deterministically, aggregating each recipe's ingredients at ingestion. The architectural difference lies in what the knowledge source is for. In \textsc{DietQA} it is the answer space: nothing can be recommended that the ingestion pipeline has not already collected, classified and -- where an ingredient matched ambiguously -- put before a human reviewer, so every extension of the data is another pass through it, and the authors note that a new locale calls for a corpus, a nutrient database and a diet taxonomy of its own. The cost surfaces in their own results, where recipe coverage falls sharply as dietary constraints accumulate and ingredient substitution has to be introduced to recover queries that would otherwise go unanswered.

In the architecture defended here the knowledge base is not the answer space but the ingredient vocabulary: the model composes the dish, and the knowledge base has only to resolve the foods it names. Nothing about a food is modelled beyond its identifying text and its nutrients, so another source, or a larger export of one already in use, is more rows behind the same lookup tool -- no schema to extend, no taxonomy to curate, no recipe to write before it can be served. The gap this thesis addresses is what that arrangement makes testable: a food knowledge base the model can be required to query for every ingredient instead of guessing, a deterministic totaller the model can call while the plan is still being composed and a critic can use to ground its review, and an ablation study isolating how much each mechanism -- separately and together -- contributes to energy and macronutrient accuracy.
